% Einheit 4 von 4 – Gleichungen aufstellen (bank/lineare-gleichungen/e4.jsonl, zusammenbau v0.1)
%% TODO Zweigzeile Teil 1: was hier gelernt wird (Satz in Schülersprache aus dem Einheitstitel) – Einheit 4
\einheitenkopf[e4]{Einheit 4 von 4 · Gleichungen aufstellen}
\zweigzeile{ab Kl. 6, je nach Buch bis Kl. 8 (am Gymnasium ab Kl. 7) · P10}
\setcounter{aufgabe}{18}

%% TODO Titel: Ich-kann-Satz fehlt in der Bank (Platzhalter: Kettenname) – Nr. 19
%% TODO Anweisung: ein Satz über der ersten Teilaufgabe fehlt in der Bank – Nr. 19
\begin{aufgabe}{Aufstellen}
%% TODO \rechenplatz{n}: Schrittzahl des Grundfalls fehlt in der Bank (nur bei fester Schreibform, 2.3 b)
\begin{teile}
\teil „Eine Zahl plus 9 ergibt 23." – welche Gleichung passt? ($x$: die gesuchte Zahl) \\ \kreuz{$9x = 23$} \\ \kreuz{$x + 9 = 23$} \\ \kreuz{$x - 9 = 23$}
\teil Ich denke mir eine Zahl. Addiere ich 18, erhalte ich 45. Gleichung und Zahl? ($x$: die gesuchte Zahl) x = \leerfeld
\teil Das Dreifache einer Zahl, vermehrt um 8, ergibt 44. Gleichung und Zahl? ($x$: die gesuchte Zahl) x = \leerfeld
\teil Mia ist 4 Jahre älter als ihr Bruder. Zusammen sind sie 28 Jahre alt. Wie alt ist der Bruder? ($x$: Alter des Bruders in Jahren) \leerfeld[Jahre]
\teil Ein Rechteck hat den Umfang 46 cm, die Seite $a$ ist 8 cm lang. Wie lang ist die Seite $b$? Schreibe gegeben, gesucht, Formel, Rechnung. b = \leerfeld[cm]
\teil Ein Freibad hatte am Samstag 380 Gäste, der Eintritt kostet 5 € pro Person. Der Betreiber sagt: „Mit $x$ Gästen mehr hätten wir 2\,400 € eingenommen." Stelle eine Gleichung auf und berechne, wie viele Gäste mehr es hätten sein müssen. \hfill (P10 2020 OS) \\ x = \leerfeld
\end{teile}
\end{aufgabe}

%% TODO Titel: Ich-kann-Satz fehlt in der Bank (Platzhalter: Kettenname) – Nr. 20
%% TODO Anweisung: ein Satz über der ersten Teilaufgabe fehlt in der Bank – Nr. 20
\begin{aufgabe}{Formel umstellen}
\begin{teile}
\teil Flächeninhalt eines Rechtecks: $A = a \cdot b$ – nach $b$ umstellen. b = \leerfeld
\end{teile}
\end{aufgabe}

%% TODO Titel: Ich-kann-Satz fehlt in der Bank (Platzhalter: Kettenname) – Nr. 21
%% TODO Anweisung: ein Satz über der ersten Teilaufgabe fehlt in der Bank – Nr. 21
\begin{aufgabe}{Aufstellen – Fehler finden, Begründen, Darstellungswechsel, Anwendung}
\begin{teile}
\teil „Eine Zahl, um 6 vermindert, ergibt 10." Lisa schreibt $6 - x = 10$. Finde den Fehler.
\teil „Das Vierfache der Summe aus einer Zahl und 3 ist 36." Tom schreibt $4x + 3 = 36$. Finde den Fehler.
\teil Für einen Ausflug wird berechnet, wie viele Busse nötig sind; die Gleichung liefert $x = 3{,}5$. Wie viele Busse werden bestellt? Begründe.
\teil $x + 12 = 30$ – Zahlenrätsel dazu, und die Zahl?
\teil Eine Taxifahrt kostet 4 € Grundpreis und 2,20 € je Kilometer. Frau Kaya zahlt 37 €. Wie lang war die Fahrt? ($x$: Strecke in km) \leerfeld[km]
\end{teile}
\end{aufgabe}
